
\documentclass{../concepts/prepareprotokol} %všechny balíčky pro protokol + titulní strana
\usepackage{graphicx}
\usepackage{newfloat}
\usepackage{caption}
\usepackage{booktabs}
\usepackage{amsmath}
\usepackage{newfloat}
\usepackage{caption}
\usepackage{siunitx}
\sisetup{inter-unit-product =$\cdot$}
\sisetup{separate-uncertainty=true}
\usepackage{subcaption}

\praktikum{II} %číslo praktika (I, II nebo III)
\autor{David Němec} %jméno autora
\datum{24.\,11.\,2025} %datum měření
\cislo{5} %číslo úlohy
\nazev{Měření osciloskopem} %název úlohy

\DeclareFloatingEnvironment[fileext=los, placement={htbp}, name=Obrázek]{schema}

\newcommand{\nobars}{
Chybové úsečky nebyly pro svou malou velikost pro přehlednost vykresleny.}


\begin{document}
\renewcommand{\figurename}{Graf}
\renewcommand{\refname}{Literatura}

\maketitle

% ========================== PRACOVNÍ ÚKOL ============================
\section{Pracovní úkol}
\begin{enumerate}
    \item Pomocí osciloskopu změřte špičkovou hodnotu napětí na svorkách sekundárního vinutí převodního transformátoru a porovnejte ji s hodnotou naměřenou na střídavém rozsahu digitálního voltmetru.
    \item Sledujte činnost jednocestného usměrňovače s křemíkovými diodami KY711
    \item Při maximální hodnotě zatěžovacího odporu \qty{10}{\kilo\ohm} sledujte závislost stejnosměrného napětí na filtrační kapacitě C v intervalu 0–10 \unit{\micro\farad}. Hodnotu usměrněného napětí při C = \qty{10}{\micro\farad} srovnejte se špičkovou hodnotou pulzního průběhu.
    \item Při filtrační kapacitě C = \qty{10}{\micro\farad} změřte závislost střídavé složky usměrněného napětí na odebíraném proudu. Proud upravujte nastavením odporu Rz. U jednocestného usměrňovače měřte do proudu \qty{0.6}{\milli\ampere}.
    \item Naměřené závislosti zpracujte graficky přímo v praktiku.
    \item Na osciloskopu zobrazte voltampérovou charakteristiku vakuové diody podle schématu připojeného k úloze. Orientačně načrtněte pozorovanou charakteristiku a vyznačte měřítka na osách. Odhadněte napětí na diodě při proudu 20 mA v propustném směru.
\end{enumerate}

% ============================= TEORIE ================================
\section{Teorie}
\subsection{Střední a efektivní hodnoty střídavého napětí}

Střední hodnota střídavého napětí $\overline{U}$ s periodickým průběhem je definována jako \cite{studijni}
\begin{equation} \label{eq:stredni}
    \overline{U} = \frac{1}{T} \int_0^T u(t) dt \, ,
\end{equation}
kde $t$ je čas, $T$ perioda a $u(t)$ je okamžitá hodnota napětí. Přístroje měřící na stejnosměrném rozsahu ukazují právě tuto střední hodnotu napětí $\bar{U}$. Zobrazovaná hodnota bude časově nezávislá, pokud intergrační doba přístroje bude menší než perioda střídavého napětí. To je u napětí odebíraného ze sítě ($f = \qty{50}{\hertz}, T = \frac{1}{f} = \qty{20}{\milli\second}$) splněno. \cite{studijni}

Efektivní hodnota střídavého napětí $U_{eff}$ je naopak definována vztahem \cite{studijni}
\begin{equation} \label{eq:efektivni}
    U_{eff}^2 = \frac{1}{T} \int_0^T u^2(t) dt \, .
 \end{equation}

V případě harmonického průběhu napětí (tj. $U = U_0 \sin(\omega t + \varphi_0)$ se špičkovou hodnotou $U_0$ a úhlovou frekvencí $\omega$) platí pro $U_{eff}$ jednoduchý vztah \cite{studijni}
\begin{equation} \label{eq:stredkueff}
    U_{eff} = \frac{U_0}{\sqrt{2}} \, .
\end{equation}

\subsection{Filtrace a usměrňování napětí}

Schéma zapojení jednocestného usměrňovače s filtrem napětí je nakresleno na obrázku \ref{sch:filtr} a vysvětleno v kapitole \ref{filtrace}. Výsledný průběh okamžité hodnoty napětí na čase je na obrázku \ref{sch:prubehy}.

\begin{schema}
\centering
\begin{subschema}{0.49\linewidth}
    \centering
    \includegraphics[width=\linewidth]{uloha5/udiod.png}
    \caption{}
    \label{sch:bezfiltru}
\end{subschema}
\hfill
\begin{subschema}{0.49\linewidth}
    \centering
    \includegraphics[width=\linewidth]{uloha5/udiodc.png}
    \caption{}
    \label{sch:sfiltrem}
\end{subschema}

\caption{Průběhy napětí usměrněného jednocestným usměrňovačem bez filtru (a) a s filtrem (b). \cite{studijni}}
\label{sch:prubehy}
\end{schema}

Usměrňovač propustí jen tu část střídavého napětí, která odpovídá proudu procházejícím diodou v propustném směru. V době, kdy dioda propouští proud se kondenzátor nabíjí přes odpor $R_{ochr}$ až na špičkovou hodnotu $U_0$. V době, kdy dioda proud nepropouští, se kondenzátor vybíjí přes odpor $R_z$ podle rovnice
\begin{equation} \label{eq:vybijeni}
    u = U_0 e^{-\frac{t}{R_z C}} \, .
\end{equation}
Za předpokladu, že $R_{ochr} \ll R_z$ a že kondenzátor se nabije za zanedbatelně krátkou dobu, dostaneme dosazením (\ref{eq:vybijeni}) do (\ref{eq:stredni}) vztah pro střední hodnotu napětí
\begin{equation} \label{eq:filtr}
    \overline{U} = \frac{U_0 R_z C}{T} \left( 1-e^{-\frac{T}{R_z C}} \right) \, .
\end{equation}
Kondenzátor se vybíjí až do příchodu následujícího pulzu v čase $t=T$. Pokud je poměr $\frac{T}{R_z C} \ll 1$, pak pro špičkovou hodnotu střídavé složky usměrněného napětí (tj. rozdíl minimální a maximální hodnoty napětí) $\Delta U$ platí (upraveno z \cite{studijni})
\begin{equation} \label{eq:deltau}
    \Delta U = U_0 - u(T) \approx \frac{U_0 T}{R_z C} \, .
\end{equation}

Stejnosměrný proud $I_{ss}$ zátěží $R_z$ je dán Ohmovým zákonem \cite{studijni}
\begin{equation} \label{eq:ohm}
    I_{ss} = \frac{U_{ss}}{R_z} \, ,
\end{equation}
kde $U_{ss}$ je stenosměrné napětí na zátěži. Za předpokladu $\frac{T}{R_z C} \ll 1$ je $U_{ss} \approx U_0$ a vztah (\ref{eq:deltau} lze přepsat na
\begin{equation} \label{deltauss}
    \Delta U = I_{ss} \frac{T}{C} \, .
\end{equation}



\subsection{Metoda měření}

\subsubsection{Filtrace jednocestně usměrněného napětí} \label{filtrace}

Jednocestný usměrňovač je na obrázku \ref{sch:filtr} vyznačen čárkovanou čarou. Skládá se z diody, která 

\begin{schema}[!htbp]
    \centering
    \includegraphics[width=0.5\linewidth]{uloha5/filtr.png}
    \caption{Schéma zapojení pro měření jednocestného usměrňovače \cite{studijni}}
    \label{sch:filtr}
\end{schema}

\subsubsection{Voltampérová charakteristika diody} \label{va}

\begin{figure}[!htbp]
    \centering
    \includegraphics[width=0.5\linewidth]{uloha5/va.png}
    \caption{Schéma zapojení pro zobrazení V-A charakteristiky diody \cite{studijni}}
    \label{sch:va}
\end{figure}



\subsection{Měřicí přístroje a jejich nejistoty}
\begin{enumerate}
    \item Multimetr Fluke 175:
    
    Multimetrem byla nejprve ověřována činnost osciloskopu a následně jím byla měřena střední hodnota usměrněného napětí (využitím stejnosměrného rozsahu multimetru). Nejistota měřené hodnoty stejnosměrného napětí je z přiložených specifikací přístroje 0.15\,\% z měřené hodnoty + 2 digity a nejistota stejnosměrného proudu je 1\,\% z měřené hodnoty + 3 digity.

    \item Odporová dekáda Metra XL6:

    Odporová dekáda představovala zátěžový odpor pro studium činnosti jednocestného usměrňovače. Nejistota nastaveného odporu je ze specifikací přístroje 0.1\,\% hodnoty.

    \item Kapacitní dekáda ELC DC05:

    Kapacitní dekáda představovala filtrující kapacitu pro studium činnosti jednocestného usměrňovače. Ve druhém úkolu pak sloužila ke kompenzaci parazitní indukčnosti odporové dekády. Nejistota nastavené kapacity je ze specifikací přístroje 1\,\% hodnoty.

    \item Osciloskop GW Instek GOS-620:

    Osciloskopem byla měřena střídavá složka usměrněného napětí a následně posloužil k zobrazení voltampérové charakteristiky vakuové diody. Jeho nejistota byla odhadnuta jako velikost nejmenšího dílku vykreslené mřížky, tedy závisela na nastaveném měřítku.
\end{enumerate}


% ======================== VÝSLEDKY MĚŘENÍ ===========================
\section{Výsledky měření}

\subsection{Špičková a efektivní hodnota napětí}


\begin{figure}
    \centering
    \includegraphics[width=0.5\linewidth]{uloha5/GRAF1.png}
    \caption{Závislost kapacity deskového kondenzátoru na napětí.}
    \label{fig:osciloskop}
\end{figure}


% ============================ DISKUZE ===============================
\section{Diskuze}


% ============================= ZÁVĚR ================================
\section{Závěr}

Špičková hodnota střídavého napětí na výstupních svorkách transformátoru byla změřena osciloskopem jako $U_0 = \qty{10\pm1}{\volt}$. Digitálním voltmetrem byla při tomto nastavení změřena efektivní hodnota střídavého napětí $U_{eff} = \qty{6.72\pm0.10}{V}$.




% =========================== LITERATURA =============================
\begin{thebibliography}{9}

\bibitem{studijni} Kolektiv ZFP KVOF MFF UK: Měření dielektrických vlastností materiálů [online]. [cit. 26.10.2025], dostupné z https://physics.mff.cuni.cz/vyuka/zfp/\_media/zadani/texty/txt\_227.pdf

\bibitem{pokyny} Kolektiv ZFP KVOF MFF UK: Měření odporu [online]. [cit. 26.10.2025], dostupné \\
z https://physics.mff.cuni.cz/vyuka/zfp/\_media/zadani/texty/txt\_202.pdf

\bibitem{englich} J. Englich: Úvod do praktické fyziky I: Zpracování výsledků měření. 1. vyd. Praha: Matfyzpress, 2006

\bibitem{mikulcak} Mikulčák a kolektiv: Matematické, fyzikální a chemické tabulky pro střední školy: relativní permitivita, Praha: Státní pedagogické nakladatelství, n. p., 1988

\end{thebibliography}
\end{document}